\documentclass[11pt,a4paper]{article}

% ============================================================================
% Sensory Tracer Science (STS): A Proposal for a New Subfield
% arXiv-ready manuscript (theory + in-silico framework)
%
% NOTE TO SUBMITTER: This source compiles with pdflatex + bibtex:
%   pdflatex main && bibtex main && pdflatex main && pdflatex main
% (or simply: latexmk -pdf main.tex)
% Update the \author block and ORCID/affiliation before submission.
% ============================================================================

\usepackage[utf8]{inputenc}
\usepackage[T1]{fontenc}
\usepackage{amsmath,amssymb,amsthm}
\usepackage{bm}
\usepackage{graphicx}
\usepackage{booktabs}
\usepackage{geometry}
\usepackage{hyperref}
\usepackage[numbers,sort&compress]{natbib}
\usepackage{xcolor}
\usepackage{enumitem}

\geometry{margin=1in}
\hypersetup{
  colorlinks=true,
  linkcolor=blue!50!black,
  citecolor=blue!50!black,
  urlcolor=blue!50!black
}

\newtheorem{axiom}{Axiom}[section]
\newtheorem{definition}{Definition}[section]

\title{\textbf{Sensory Tracer Science (STS): A Conservation-Law\\
Framework for Auditable Signal Propagation in\\
Biological and Computational Media}}

% ----------------------------------------------------------------------------
% TODO (submitter): replace with your real name, affiliation, email, ORCID.
% ----------------------------------------------------------------------------
\author{%
  Kevin Kull\thanks{Corresponding author: \texttt{kevinkull.kk@gmail.com}.
  Affiliation and ORCID to be completed by the submitting author.}\\
  \small Sensory Tracer Science Project\\
  \small \url{https://github.com/Kuonirad/Sensory-Tracer-Science-STS-Proposal-for-a-New-Subfield}
}

\date{\today}

\begin{document}
\maketitle

\begin{abstract}
We propose \emph{Sensory Tracer Science} (STS), a candidate subfield that studies
the \emph{energy-conserving}, \emph{information-preserving}, and
\emph{causality-respecting} propagation of sensory signals through distributed
media. STS introduces no new physics. Instead it imposes a stricter
\emph{accounting} discipline: every traced signal must close three independent
ledgers---an energy balance grounded in Landauer's principle and the first law of
thermodynamics, an information balance grounded in Shannon theory, and a causality
constraint grounded in the finite speed of propagation in a medium. We formalize
the framework as five axioms ($\mathrm{A1}$--$\mathrm{A5}$) reduced from
established results (Landauer, Shannon, Heisenberg, special relativity, and
cellular bioenergetics), derive a governing set of continuity and wave-propagation
equations, and define a \emph{triple-audit} validation protocol with explicit
numerical tolerances. We present a reference open-source implementation that
realizes the framework across three modalities in simulation---fiber-optic
(Brillouin) tracers, biocompatible neural tracers, and quantum-enhanced
tracers---and that subjects each to the same three audits. We then show that the
\emph{same} audit engine applies, unchanged, to a second domain: by treating the
activation flow of a reasoning/inference stack as ``synthetic neural tissue,'' STS
yields an end-to-end energy/information/causality account of how a signal
propagates through a computational pipeline, making auditable attribution a
first-class primitive. All results reported here are theoretical and in-silico;
we make no wet-lab validation claims, and we are explicit about which quantities
are derived from first principles versus parameterized from the literature. We
argue that the principal contribution of STS is not any single tracer but a
falsifiable, conservation-first methodology---one we apply reflexively to the
project's own software and documentation.
\end{abstract}

\noindent\textbf{Keywords:} sensory tracers; information thermodynamics;
Landauer's principle; conservation laws; biophysics; neural tracing; quantum
sensing; interpretability; auditable attribution.

\section{Introduction}
\label{sec:intro}

Tracing a signal through a medium is a recurring problem across the sciences: a
fluorophore through neural tissue, an acoustic phonon through an optical fiber, an
entangled photon through a detector array, or---increasingly---an activation
through the layers of a computational inference stack. Each community has its own
validation idiom: histology for neural tracers, tomography for quantum sensing,
calibration for distributed sensor networks. What these idioms share, and what no
single community enforces uniformly, is a commitment to the conservation laws that
ultimately bound \emph{any} physical signal: energy cannot be created, information
cannot be detected where it was never injected, and nothing propagates faster than
the medium allows.

We propose \emph{Sensory Tracer Science} (STS) as a subfield organized entirely
around this commitment. The central thesis is deliberately conservative:

\begin{quote}
\itshape STS introduces no new physics. It introduces a stricter accounting of
existing physics, applied uniformly to every traced signal regardless of the
medium it traverses.
\end{quote}

This ``zero-new-physics'' stance is a feature, not a limitation. By refusing to
add laws and instead demanding that every implementation \emph{close its books}
against three independent ledgers, STS produces claims that are falsifiable in the
strict sense: an implementation either passes its energy, information, and
causality audits within stated tolerances, or it does not. The unit of validation
is the audit, not the anecdote.

\paragraph{Contributions.}
\begin{enumerate}[leftmargin=1.4em]
  \item We axiomatize STS as five constraints reduced from established physics
        (Section~\ref{sec:axioms}) and a small set of governing equations
        (Section~\ref{sec:equations}).
  \item We define a \emph{triple-audit} validation protocol with explicit
        numerical tolerances and a binary pass/fail outcome
        (Section~\ref{sec:audit}).
  \item We describe a reference implementation that applies the protocol to three
        physical tracer modalities in simulation
        (Section~\ref{sec:domains}).
  \item We show the same audit engine extends, without modification, to
        computational pipelines by treating activation flow as synthetic neural
        tissue---turning end-to-end attribution into an auditable primitive
        (Section~\ref{sec:bridge}).
  \item We discuss the project's reflexive application of the same
        conservation/audit discipline to its own software and prose
        (Section~\ref{sec:meta}), and we state the framework's limitations
        plainly (Section~\ref{sec:limitations}).
\end{enumerate}

\paragraph{Scope and honesty statement.}
Everything reported in this manuscript is theoretical or obtained from
deterministic in-silico simulation of the framework's own equations. We make
\emph{no} claim of wet-lab or clinical validation. Where a quantity is derived
from first principles we say so; where it is parameterized from the experimental
literature (e.g.\ diffusion coefficients, ATP free energy, two-photon cross
sections) we label it as such and cite the source. Distinguishing these two
categories is itself part of the STS accounting discipline.

\section{Foundational Axioms}
\label{sec:axioms}

STS rests on five axioms. Each is a constraint \emph{reduced from} an established
result; none asserts a new law.

\begin{axiom}[Information--energy conservation]\label{ax:a1}
Information and energy are distinct but jointly bounded by conservation laws.
Erasing one bit dissipates at least $k_B T \ln 2$ joules (Landauer's
principle~\citep{landauer1961irreversibility,bennett1982thermodynamics}), and
every continuous symmetry corresponds to a conserved quantity. In a closed
system, energy used to encode information is neither created nor destroyed, and
information loss manifests as entropy increase and heat.
\end{axiom}

\begin{axiom}[Speed limit in a medium]\label{ax:a2}
No tracer propagates faster than $v \le c/n$, where $c$ is the vacuum speed of
light and $n\ge 1$ the medium's refractive index. This is the causality
constraint.
\end{axiom}

\begin{axiom}[Second law]\label{ax:a3}
Total entropy production is non-negative, $\Delta S_\mathrm{total}\ge 0$; any local
decrease must be compensated globally. No tracer is perfectly efficient and heat
production is unavoidable.
\end{axiom}

\begin{axiom}[Measurement uncertainty]\label{ax:a4}
Measurement injects noise bounded below by the Heisenberg
limit~\citep{heisenberg1927uncertainty}: conjugate-variable products satisfy
$\Delta x\,\Delta p \ge \hbar/2$. This sets the ultimate sensitivity floor.
\end{axiom}

\begin{axiom}[Biological energy budget]\label{ax:a5}
In biological media, sustainable signal processing is bounded by metabolic supply.
ATP hydrolysis provides on the order of $|\Delta G_\mathrm{ATP}|\approx
57\,\mathrm{kJ\,mol^{-1}}$ (pH 7.4, $1\,\mathrm{mM\;Mg^{2+}}$,
$37^{\circ}\mathrm{C}$), capping the rate of information processing in living
tissue.
\end{axiom}

These axioms are governed by five meta-axioms enforcing internal consistency:
every axiom must be derivable from physical observables; no two may contradict
under any conditions; any violation must be traceable to a single axiom; the
framework must be invariant under reference-frame transformations; and---most
importantly---STS introduces no new physics, only new organization of existing
physics.

\section{Governing Equations}
\label{sec:equations}

\paragraph{Conservation of sensory information.}
The total traced information is modeled as a conserved functional of sensor
density $\rho_\mathrm{sensor}$ and local tracer energy $E_\mathrm{tracer}$:
\begin{equation}
I_\mathrm{total}
= \iiint \rho_\mathrm{sensor}(\bm{r},t)\,
  \log_2\!\left(1 + \frac{E_\mathrm{tracer}(\bm{r},t)}{k_B T}\right)
  \, d^3r\, dt = \text{const},
\label{eq:info}
\end{equation}
where the logarithmic term is the Shannon channel capacity of a thermally limited
link~\citep{shannon1948mathematical}.

\paragraph{Tracer energy continuity.}
Energy density obeys a continuity equation with non-negative dissipation:
\begin{equation}
\frac{\partial E_\mathrm{tracer}}{\partial t}
+ \nabla\cdot \bm{J}_\mathrm{tracer}
= -\,P_\mathrm{dissipation} + P_\mathrm{source},
\qquad P_\mathrm{dissipation}\ge 0,
\label{eq:continuity}
\end{equation}
which enforces Axioms~\ref{ax:a1} and~\ref{ax:a3} pointwise.

\paragraph{Damped wave propagation.}
Tracer fields propagate as damped, driven waves
\begin{equation}
\frac{\partial^2 \psi}{\partial t^2}
= v^2 \nabla^2 \psi - \gamma\,\frac{\partial \psi}{\partial t}
+ S_\mathrm{sensor}(\bm{r},t),
\qquad v \le c/n,
\label{eq:wave}
\end{equation}
with attenuation coefficient $\gamma\ge 0$ and the speed bound from
Axiom~\ref{ax:a2}.

\paragraph{Representative closed-form predictions.}
From the equations above and the axioms, several quantities follow directly and
serve as testable predictions:
\begin{align}
E_\mathrm{min} &= k_B T \ln 2 &&\text{(minimum energy per bit)},\\
L_\mathrm{max} &= \frac{v}{\gamma}\,\ln\!\frac{E_0}{k_B T}
  &&\text{(max.\ range before signal $\to$ thermal floor)},\\
T_2 &\le \frac{\hbar}{2 k_B T} &&\text{(thermal coherence-time bound)}.
\end{align}
At $T=300\,\mathrm{K}$, $E_\mathrm{min}\approx 2.87\times10^{-21}\,\mathrm{J\,bit^{-1}}$;
the coherence bound gives $T_2 \lesssim 1.3\times10^{-14}\,\mathrm{s}$. These are
order-of-magnitude consistency checks, not engineering specifications.

\section{The Triple-Audit Validation Protocol}
\label{sec:audit}

The operational core of STS is a validation protocol that every implementation
must pass. It comprises three independent audits, each with an explicit tolerance,
and yields a binary outcome: failure of \emph{any} audit invalidates the system.

\begin{enumerate}[leftmargin=1.6em]
  \item \textbf{Energy audit.}
  \begin{equation}
  \bigl| E_\mathrm{in} - (E_\mathrm{out} + E_\mathrm{dissipated}) \bigr|
  \;<\; \varepsilon_E\, E_\mathrm{in},
  \end{equation}
  with $\varepsilon_E = 10^{-12}$ (picojoule-per-joule, set by numerical
  precision). Separately, dissipation is checked against the Landauer floor
  $k_B T \ln 2$ per erased bit as a physical-plausibility flag.

  \item \textbf{Information balance.}
  \begin{equation}
  \bigl| I_\mathrm{injected} - (I_\mathrm{detected} + I_\mathrm{lost}) \bigr|
  \;<\; \varepsilon_I\, I_\mathrm{injected},
  \qquad \varepsilon_I = 10^{-2}.
  \end{equation}
  Shannon content and step-wise mutual information are computed along the trace.

  \item \textbf{Causality check.}
  \begin{equation}
  v_\mathrm{signal} \le v_\mathrm{medium}, \qquad \text{tolerance } = 0
  \end{equation}
  (zero tolerance for faster-than-medium propagation), augmented by a
  time-of-flight check.
\end{enumerate}

The reference implementation realizes these as three auditor classes
(\texttt{EnergyAuditor}, \texttt{InformationAuditor}, \texttt{CausalityAuditor})
composed by a top-level \texttt{STSValidator}. Constants are pinned to CODATA
recommended values~\citep{codata2018}.

\section{Implementation Domains (In-Silico)}
\label{sec:domains}

The reference framework instantiates the protocol across three modalities. All
results below are from deterministic simulation of the framework's equations under
the stated constraints; none are wet-lab measurements.

\subsection{Fiber-optic (Brillouin) tracers}
Acousto-optic tracing via Brillouin
scattering~\citep{brillouin1922scattering}, constrained to
$<1\,\mathrm{nJ}$ input to avoid nonlinear damage, with silica refractive index
$n=1.46$ and telecom-band attenuation. The audit verifies energy conservation in
the optical--acoustic interaction.

\subsection{Biocompatible neural tracers}
Diffusion--advection with metabolic clearance, parameterized from the literature:
free $\mathrm{Ca^{2+}}$ diffusion ($D\approx 2.2\times10^{-10}\,\mathrm{m^2\,s^{-1}}$,
$37^{\circ}\mathrm{C}$)~\citep{einstein1905motion,fick1855diffusion}, cytoplasmic
viscosity, blood--brain-barrier clearance, and ATP-bounded energy budgets
(Axiom~\ref{ax:a5}). Modeled constraints include tracer concentration
$<1\,\mu\mathrm{M}$, ATP depletion rate $<0.1\,\mathrm{mM\,s^{-1}}$, and
phototoxic dose $<50\,\mathrm{J\,mm^{-3}}$. Supporting electrochemistry uses the
Nernst~\citep{nernst1888theory} and Debye--H\"uckel~\citep{debye1923theory}
relations, with membrane biophysics following
Hodgkin--Huxley~\citep{hodgkin1952quantitative} and Hille~\citep{hille2001ionic}.

\subsection{Quantum-enhanced tracers}
Entangled-photon tracing constrained by photon flux
($<10^9\,\mathrm{photons\,pixel^{-1}\,s^{-1}}$), antibunching
($g^{(2)}(0)<0.1$), and the thermal coherence bound of Axiom~\ref{ax:a4}. The
audit enforces the Heisenberg-limited noise floor.

\paragraph{Uncertainty handling.}
Where literature parameters carry uncertainty (e.g.\ logBB, $k_\mathrm{off}$,
quantum yield, photobleaching rate), the framework propagates Gaussian
uncertainties by Monte Carlo ($N=10^4$) and reports the resulting spread rather
than point estimates.

\section{From Tissue to Reasoning Stacks: Auditable Attribution as a Primitive}
\label{sec:bridge}

A computational inference pipeline shares the STS primitives: a kernel transforms
an activation state, that transformation costs energy (Landauer-bounded compute),
it preserves or loses information, and steps form a directed acyclic graph (DAG)
of strictly ordered causality. STS exploits this structurally rather than
metaphorically: it treats an activation flow as \emph{synthetic neural tissue} and
runs the \emph{same} triple audit over the stack.

Concretely, given a sequence of kernels $K_1,\dots,K_m$ (duck-typed callables
$\mathbb{R}^d\!\to\!\mathbb{R}^d$) and an initial activation $\bm{a}_0$, an
attributor produces an enriched trace reporting:
\begin{itemize}[leftmargin=1.4em]
  \item \textbf{energy efficiency} $E_\mathrm{out}/E_\mathrm{in}$ over the chain,
        with propagation accounting constructed to satisfy
        $E_\mathrm{in}=E_\mathrm{out}+E_\mathrm{dissipated}$ exactly so the energy
        audit passes by construction, while the Landauer floor is flagged
        separately;
  \item \textbf{information fidelity} $I_\mathrm{detected}/I_\mathrm{injected}$ via
        Shannon content and step-wise mutual information;
  \item \textbf{Landauer efficiency}: dissipation relative to the $k_B T \ln 2$
        bound, a physical-realism flag;
  \item \textbf{causality}: DAG acyclicity plus a normalized medium-speed budget
        per transition.
\end{itemize}

The payoff is that \emph{end-to-end attribution becomes a first-class, auditable
primitive}: any reasoning loop can be wrapped to obtain a conservation-checked
account of how a signal propagates through it. The same three ledgers that
validate a biological tracer validate a computational one. This is the dual-use
bridge at the heart of the proposal---one engine, two domains, no new physics. We
emphasize that the energy figures here are an \emph{accounting} abstraction over
activation norms, not a hardware power measurement; the contribution is the
auditable bookkeeping, not a claim about wall-clock joules.

\section{Reflexive Integrity: Auditing the Project Itself}
\label{sec:meta}

STS applies its own discipline reflexively. The reference project ships
self-validation scripts that re-derive constants and check axiom consistency at
import time, machine-checked audit artifacts, and readability tooling that holds
the documentation to a measured standard. The intent is methodological
consistency: a framework whose thesis is ``close your books'' should close its
own. We regard this self-referential accounting as part of the contribution, not
decoration---it makes the project's claims about itself as falsifiable as its
claims about tracers.

\section{Comparison with Existing Fields}
\label{sec:comparison}

\begin{table}[t]
\centering
\small
\begin{tabular}{@{}lllll@{}}
\toprule
\textbf{Feature} & \textbf{STS} & \textbf{Neural tracing} & \textbf{Quantum sensing} & \textbf{Distributed sensors}\\
\midrule
Conservation law   & derived & empirical & quantum & engineering\\
Energy bound       & $k_B T\ln 2$ & ATP-limited & $\hbar\omega$ & battery life\\
Causality          & strict & approximate & strict & approximate\\
Information metric  & $I_\mathrm{total}=$const & connectivity map & Fisher info & SNR\\
Validation         & triple audit & histology & tomography & calibration\\
\bottomrule
\end{tabular}
\caption{STS contrasted with adjacent fields. STS's distinguishing feature is a
uniform, derived, triple-audit validation applied across media.}
\label{tab:compare}
\end{table}

Table~\ref{tab:compare} situates STS. It is not a competitor to these fields but a
cross-cutting methodology: any of them could adopt the triple audit without
abandoning its native instrumentation.

\section{Limitations and Threats to Validity}
\label{sec:limitations}

We state limitations plainly, consistent with the framework's accounting ethos.
\begin{enumerate}[leftmargin=1.6em]
  \item \textbf{No wet-lab validation.} All quantitative results are in-silico.
        The biological constraints are parameterized from the literature, not
        measured in this work. Claims of biocompatibility or clinical readiness
        are \emph{not} made here.
  \item \textbf{Audit-by-construction.} In the computational bridge, the energy
        ledger is satisfied by construction; its value is bookkeeping
        transparency, not an independent measurement of compute energy.
  \item \textbf{Model abstraction.} Equations~\eqref{eq:info}--\eqref{eq:wave} are
        deliberately reduced; real media exhibit nonlinearity, anisotropy, and
        active transport beyond the modeled terms.
  \item \textbf{Scope of novelty.} STS contributes organization and methodology,
        not new physical law; readers seeking new dynamics will not find them
        here, by design.
\end{enumerate}

\section{Conclusion}
\label{sec:conclusion}

We have proposed Sensory Tracer Science as a conservation-first subfield whose
unit of validation is a triple audit---energy, information, causality---applied
uniformly to any traced signal, biological or computational. The framework adds no
physics; it adds discipline. Its reference implementation demonstrates the
protocol in simulation across three physical modalities and, using the identical
audit engine, across computational inference stacks, making auditable end-to-end
attribution a first-class primitive. We view the methodology, applied reflexively
to the project's own artifacts, as the durable contribution. The natural next step
is precisely the one STS is built to invite: an adversarial, wet-lab test that
tries to make a real tracer \emph{fail} its books.

\section*{Code and Data Availability}
The reference implementation, self-validation scripts, and this manuscript source
are available at
\url{https://github.com/Kuonirad/Sensory-Tracer-Science-STS-Proposal-for-a-New-Subfield}
under the MIT license. No new experimental datasets were generated.

\section*{Acknowledgments}
We thank the open-source scientific-computing community whose libraries and
constants (CODATA) underpin the reference implementation.

\bibliographystyle{unsrtnat}
\bibliography{references}

\end{document}
